
\documentclass[11pt]{article}
\usepackage{enumitem,amssymb}
\newlist{todolist}{itemize}{2}
\setlist[todolist]{label=$\square$}

\usepackage{graphicx}

      \textwidth 15cm
      \textheight 22cm
      \parindent 10pt
      \oddsidemargin 0.85cm
      \evensidemargin 0.37cm

\usepackage[utf8]{inputenc}
\usepackage{tcolorbox}
\usepackage{xcolor}
% to break the long urls 
\PassOptionsToPackage{hyphens}{url}
\usepackage[colorlinks=true, urlcolor=blue, pdfborder={0 0 0}, pdftex,
            pdfauthor={Animesh Trivedi},
            pdftitle={Thesis content at VU Amsterdam},
            pdfsubject={Thesis content at VU Amsterdam},
            pdfkeywords={MSc thesis, Bsc Thesis, advice, vu amsterdam, animesh trivedi}, 
            pdfproducer={Latex with hyperref},
            pdfcreator={pdflatex}]{hyperref}

% atr: This sets up sans or roman fonts 
%\renewcommand{\familydefault}{\sfdefault}


\begin{document}

\thispagestyle{empty}

\begin{center}

Vrije Universiteit Amsterdam

\vspace{1mm}

\includegraphics[height=28mm]{vu-griffioen.pdf}

\vspace{1.5cm}

{\Large Bachelor/Master Thesis}

\vspace*{1.5cm}

\rule{.9\linewidth}{.6pt}\\[0.4cm]
{\huge \bfseries Notes on CS BSc and MSc Thesis \par}
{\huge \bfseries  at VU Amsterdam \par}\vspace{0.4cm}
\rule{.9\linewidth}{.6pt}\\[1.5cm]

\vspace*{2mm}

%{\Large
{\normalsize 
\begin{tabular}{ll}
{\bf Author:} & Animesh Trivedi \\
web: & \url{https://animeshtrivedi.github.io/} \\ 
StoNet Research Group & \url{https://stonet-research.github.io/} \\ 
\end{tabular}
}


\vspace*{1.5cm}

\begin{tabular}{ll}
{\it 1st supervisor:}   & ~~N/A \\
{\it daily supervisor:} & ~~N/A \\
{\it 2nd reader:}       & ~~N/A 
\end{tabular}

\vspace*{2cm}

\textit{A thesis submitted in fulfillment of the requirements for\\ the VU Bachelor/Master of Science degree in Computer Science }

\vspace*{1cm}

\today\\[4cm] % Date

\end{center}

\newpage

% \author{Animesh Trivedi \\ a.trivedi@vu.nl \\ \url{https://animeshtrivedi.github.io/} \\ version: v1.5}
% \date{May 19th, 2021}


% \begin{document}
% 
% \maketitle
% \newpage
\tableofcontents
\newpage 
\section*{Version History}
\begin{itemize}
    \item v2.0     : Updated details on each chapter, and a new thesis checklist.
    \item v1.5     : Added references and tex template link to the Reproducibility Appendix. 
    \item pre-v1.4 : Initial versions. 
\end{itemize}

\newpage 
  \begin{todolist}
    \item List item 1 goes here.
    \item List item 2 goes here.
    \begin{todolist}
      \item Sublist item 1 goes here.
      \item Sublist item 2 goes here.
    \end{todolist}
    \item List item 3 goes here
    \item List item 4 goes here.
  \end{todolist}

\newpage 
\section*{What is this about}
Over the past years I have given a common subset of advice regarding how to write your thesis to multiple BSc and MSc students. This document is a synthesis of these advises, dos and don'ts that I have told them. In the coming section, I will add details of the kind of content you should be covering in your thesis and what I expect to see in each section of your thesis. BSc and MSc theses differ in the scope of the problem that they cover. Also, these are a broader set of advice that will also help when you are writing a research paper. 


I am happy to have input from you in this matter and let me know if I have missed something that you covered in your thesis, and/or how we can improve the instructions here. This is not meant to be a comprehensive set of instructions, and following these guidelines does not imply getting the maximum grade. 

\subsection*{General Guidelines}
\begin{itemize}
    \item Aim to produce a first complete draft 2-3 weeks before the final hand-in deadline. The earlier you give me a draft to review, more detailed feedback you will get. 
    \item Spell check and proofread your draft before you give me. Do not leave half broken sentences and empty subsections.
    \item Aim to write frequently, and ask questions early. 
    \item You will have one major and one minor revision opportunity before you hand in your final thesis draft for grading. It is unreasonable for me to read and re-read your drafts every week, hence, aim to write a first as complete draft as possible for maximum feedback. 
    \item When writing \textbf{\textit{"Always think like a compiler - have I explained and/or defined all the ideas that I am about to use?"}}. It is a responsibility of the writer to write an easy to read document, not the other way around. Do not just write for the writing sake to fill up pages. Try to make a simple, linear structure of your thesis. Build your case step by step, one new idea at a time, one new configuration or complexity. Try to reduce the "cognitive load" of your thesis as much as possible. If you want to test what does \textit{"cognitive load"} mean, then try to read your own thesis very late in the night when you are really tired or fresh in the morning when you do not remember all the details up front. Re-read what you have written after a couple of days to see if you remember what is the key message that you were trying to get across. 
    \item Always read and re-read your own thesis before you submit for a review. 
    \item Your thesis PDF is for everyone else than you and me. Think of your friend who also studies Computer Science but does not know what you work on. They are not invested in your problem, and are not talking to you over the last couple of months. Hence, think and write down all the details which are necessary for them to understand what you have done. 
\end{itemize}

\subsection*{English Usage}
Pay attention to your English usage, especially the use of a paragraph and subsection. Each paragraph is supposed to cover one logical idea. Do not ramble. Do not cover 100 different possibilities and ideas in a single paragraph. After each paragraph, consider does it explain one idea and one idea only cleanly. Do not write paragraphs which are 1 or 2 pages. For every paragraph connect it back and forward. Similar advice goes for subsections. Think, if the heading of the subsection makes sense. If I hide the title of your subsection, would a reader come up with a similar heading if she or he reads the content of your subsection. For every new idea: write what you are about to explain, explain it, and recap what you have explained. 


Do not use colloquial terms in your writing. Read papers and see how they present information and ideas. Do not use terms like "may", "could", "perhaps" - try to write as cleanly as possible without ambiguity. Always write in the present tense in active voice. Do not use frivolous adjectives and/or adverbs like "very much", "tremendous", "very fast", "high overheads". Always aim to quantify. What does "fast" or  "slow" mean, can you measure it? in what term, what context? Be as specific as possible when you report information. 

\textbf{The golden rule of writing} is (follow this for each subsection, section, and chapter) (i) Write what you are going to tell the reader and why is it important; (ii) Write and explain it to the reader; (iii) Recap what you have just explained to the reader.  

\subsection*{Mandatory Reading List}
Please read the following references carefully as they cover many of ideas which are iterated throughout this document:  
\begin{enumerate}
    \item George D. Gopen \& Judith A. Swan, The Science of Scientific Writing, \url{https://github.com/animeshtrivedi/notes/blob/master/docs/the-science-of-scientific-writing.pdf}. 
    \item Gernot Heiser, Tips and Guidance for Students Writing Papers and Reports, \url{https://www.cse.unsw.edu.au/~gernot/style-guide.html}. 
    \item Gernot Heiser, Systems Benchmarking Crimes, \url{https://www.cse.unsw.edu.au/~gernot/benchmarking-crimes.html}. 
    \item John Ousterhout, Always Measure One Level Deeper, \url{https://m-cacm.acm.org/magazines/2018/7/229031-always-measure-one-level-deeper/fulltext}. 
    \item Roy Levin, and David D. Redell,  How (and how not) to write a good systems paper, (applicable to your thesis work as well), \url{https://www.usenix.org/conferences/author-resources/how-and-how-not-write-good-systems-paper}.
    \item Inclusion and Diversity in Writing, \url{https://acmsocc.github.io/2020/inclusion_and_diversity_in_writing.html}. 
\end{enumerate}
\subsection*{In the End}
Take time to develop writing skills. It is going to be good from one day to another. It is as important as learning how to code, if not more. Write drafts and seek feedback. You can be creative in your writing once you have mastered the basics. 

Most importantly, your thesis work is probably the biggest piece of work you have done so far in your studies. Care for it. Take pride in your work and writing - your name is attached to it! 

\newpage
\Chpater{Abstract}
An abstract is a compressed or zip summary of your thesis. You can follow a simple structure:

\noindent (i) the broader societal, economic, scientific context, who are the people who might be interested in your work (2-4 sentences);

\noindent (ii) what is changing and what broad problem(s) does this change creates, has the problem always been here or something has triggered this problem, is it a specific problem or part of a bigger trend in your field (2-4 sentences); 

\noindent (iii) what specific scientific/research problems that this thesis focuses on (4-6 sentences); 

\noindent (iv) how does this thesis work tackle the problem (4-6 sentences); 

\noindent (v) Key findings, give "quantitative results". For example, \texttt{Our results show that our system decreases SQL query execution time by 50.6\% in the cloud} (1-2 sentence); 

\noindent (vi) How can I use this thesis work. For example, \texttt{The code, and data is available at https://...}  


These parts are guidelines, not strict rules (except the last one). Based on the flavor of the thesis you do, you will budget your word/sentence quotas differently. 

\newpage 

\newpage 

\newpage 

\newpage 

\newpage 
\newpage 

\newpage

\newpage 
\section{References}
The final section with references. The thesis template will generate the links in references for you. Please make sure that all references are consistent in formatting, and there are no broken links. Please do not leave for me to check and correct broken links for you. It is very annoying. Be vigilant and thorough. 

\newpage 

%\section{Students who helped with this document}
%Jonas Theis, 
\end{document}
